\begin{pastexam}[label={px:ntu-106-ie-ds-03}]{臺大工工所 106 學年度 第 3 題}
（是非；對填 (A)、錯填 (B)；答對 4 分，答錯倒扣 4 分）

Even though the minimum data of a min-heap is always on the root of the tree and it takes O(1) time to
find it, the removal of this minimum data may take O(log n) time as we need to arrange the nodes to
maintain the tree structure.
\end{pastexam}

\begin{pastsoln}
答案：(A) 對。

最小堆積（資結書第~8.1~節）的根是最小值，找最小值只要看 $A[1]$，$\Theta(1)$。 刪除最小值（資結書虛擬碼~8.2）：
把最後一個元素搬到根、堆積大小減一，再從根往下調整，每層和較小的子節點比，必要時交換。 調整最多走樹高那麼多層，
$n$ 個元素的堆積是完全二元樹，高度 $\lfloor \lg n \rfloor + 1$，所以刪除是 $O(\log n)$；搬上去的是最後一個葉子，
通常很大，常常一路換到底，最壞就是 $\Theta(\log n)$。 題目說「may take $O(\log n)$」，原因是要重新調整節點、
維持堆積的結構，敘述正確。
\end{pastsoln}
